\chapter{Anomaly Detection} \label{chap:AD}

As previously discussed in Sec. \ref{sec:SM_BSM/BSM}, SM is not a complete theory and New Physics beyond it should exist.
The usual model-dependent searches for BSM physics rely on MC simulations generated for specific BSM hypotheses to define signal-enriched regions.
They provide optimized sensitivity for the specific process although they remain blind to a broader scenario and they have not yet produced any significant discoveries.
This motivates the exploration of alternative approaches.\\
\textit{Model-independent searches} are the most common solution, as they do not rely on specific theoretical assumptions \cite{LHCOlympics}. 
They have been used for a long time among the High Energy Physics (HEP) community to derive limits on the production cross-section of BSM signals. 
Despite their name, these searches were still too dependent on the theoretical assumptions or on the simulation for the event selection and the background modelling, limiting their power.
Beyond that, their sensitivity was limited by the look-elsewhere effect \cite{LEE_HEP}.\\
Anomaly Detection (AD) addresses these constraints by identifying rare patterns without assuming a explicit signal hypotheses.

\section{Overview of Anomaly Detection}
\label{sec:ad_overview}

Anomaly Detection encompasses algorithms that identify data instances deviating significantly from nominal distributions \cite{belis2023machine, tinari2024}. 
In standard industrial contexts, anomalies manifest as out-of-sample or off-manifold events, such as fraudulent transactions or equipment failures \cite{kasieczka2021lhc, tinari2024}. 
In High Energy Physics (HEP), single events rarely provide statistical evidence of new phenomena; BSM signals typically manifest as subtle over-densities in phase space \cite{belis2023machine, kasieczka2021lhc}.

Within particle physics, ML-based AD strategies operate across different levels of supervision \cite{belis2023machine, tinari2024}. 
Weakly supervised methods, such as Classification Without Labels (CWoLa) and CATHODE, estimate localized over-densities by comparing signal-enriched regions against adjacent sidebands \cite{belis2023machine, tinari2024}. 
Conversely, fully unsupervised methods operate solely on background processes or unlabeled experimental data \cite{belis2023machine, tinari2024}. 
These algorithms isolate rare topology features in kinematic distributions, providing an unbiased tool for model-agnostic BSM searches \cite{belis2023machine, davide_proc}.

\section{Deep Learning and Neural Networks}\label{sec:dl_dnn}
Deep Learning (DL) models use stacked layers of parametrized mathematical transformations to extract abstract features from complex data inputs \cite{tinari2024, Russo_thesis}. 
A standard Deep Neural Network (DNN) processes an input vector $\mathbf{x} = (x_1, \dots, x_n)$ through hidden layers composed of individual processing units \cite{tinari2024, Russo_thesis}. 
The output $z$ of a single unit is defined as:\begin{equation}z = \sum_{i=1}^{n} w_i x_i + b\end{equation}where $w_i$ denotes trainable weights and $b$ represents the bias term \cite{tinari2024}. 
Non-linear activation functions $\sigma(z)$ map $z$ to break linear combinations across stacked layers \cite{tinari2024, Russo_thesis}:
\begin{equation}
y = \sigma(z)
\end{equation}
Common activation functions include the Rectified Linear Unit, $\text{ReLU}(z) = \max(0, z)$, and its smooth variant, the Gaussian Error Linear Unit \cite{tinari2024}:
\begin{equation}
    \text{GELU}(z) = z \cdot \Phi(z)
\end{equation}
where $\Phi(z)$ is the standard Gaussian cumulative distribution function \cite{tinari2024}. 
Network parameters update via gradient descent backpropagation, minimizing a task-specific loss function evaluating model predictions \cite{kasieczka2021lhc, tinari2024}.

\section{Graph Neural Networks}\label{sec:gnn}
Standard feedforward architectures struggle to process variable-size physical objects without imposing artificial ordering \cite{dav_ad_proc, Russo_thesis}. Hadronic jets, composed of varying counts of collimated sub-particles, naturally map to graph structures $G = (V, E)$, where $V$ denotes constituent nodes and $E$ represents interaction edges \cite{kasieczka2021lhc, davide_proc}. Node features contain kinematic variables $\mathbf{x}_i = (p_T, \eta, \phi)$, while edge features capture pairwise angular separations $\Delta R = \sqrt{(\Delta \eta)^2 + (\Delta \phi)^2}$ \cite{kasieczka2021lhc, davide_proc}.Graph Neural Networks (GNNs) preserve permutation invariance while executing local message passing \cite{kasieczka2021lhc, davide_proc}. The representation $\mathbf{h}_i^{(l)}$ of node $i$ at layer $l$ updates according to aggregated messages from its spatial neighborhood $\mathcal{N}(i)$ \cite{davide_proc}:\begin{equation}\mathbf{m}_{ij}^{(l)} = \phi_e \left( \mathbf{h}_i^{(l-1)}, \mathbf{h}j^{(l-1)}, \mathbf{e}{ij} \right)\end{equation}\begin{equation}\mathbf{h}i^{(l)} = \phi_v \left( \mathbf{h}i^{(l-1)}, \bigoplus{j \in \mathcal{N}(i)} \mathbf{m}{ij}^{(l)} \right)\end{equation}where $\phi_e$ and $\phi_v$ denote neural operations, $\mathbf{e}_{ij}$ is the edge feature, and $\bigoplus$ represents a permutation-invariant aggregation operator such as sum, mean, or max \cite{kasieczka2021lhc, davide_proc}. Graph-level outputs obtain global representations via permutation-invariant pooling functions across all nodes \cite{davide_proc}.

\section{Anomaly Detection Implementation via Deep Architectures}\label{sec:ad_implementation}
Unsupervised AD implementation utilizing DNNs and GNNs relies on latent representation learning \cite{belis2023machine, tinari2024}. Autoencoders (AEs) compress input graph representations $\mathbf{x}$ into a lower-dimensional bottleneck $\mathbf{z} = f_{\text{enc}}(\mathbf{x})$ before generating a reconstructed output $\hat{\mathbf{x}} = f_{\text{dec}}(\mathbf{z})$ \cite{belis2023machine, tinari2024}. When trained exclusively on dominant Standard Model background events, the network learns lossy compression tailored to nominal physics patterns \cite{belis2023machine, davide_proc}.Events containing BSM physics produce higher reconstruction errors due to unlearned substructure features \cite{belis2023machine, tinari2024}. An Anomaly Score ($AS$) is defined using the Mean Squared Error (MSE) metric \cite{belis2023machine, dav_ad_proc}:\begin{equation}AS = \frac{1}{N} \sum_{i=1}^{N} \left| \mathbf{x}_i - \hat{\mathbf{x}}_i \right|^2\end{equation}In ATLAS model-independent searches, cuts applied to the calculated $AS$ distribution isolate BSM candidates \cite{belis2023machine, davide_proc}. By filtering out high-probability background interactions, these architectures define signal-enriched regions directly from data, enabling model-agnostic resonance searches \cite{belis2023machine, tinari2024}.

% #################################################à
\section{What is it?}\label{sec:WhatisAD}
When applying Anomaly Detection \cite{DLforAD}\cite{LHCOlympics}, we are looking for outliers or rare events, which typically have a small deviation from already known events. In this case data is divided into two categories: "background", containing Standard Model events or noise, and "signal", including the hypothetical New Physics. All SM phenomena are classified within background without any distinction based on the origin of data. It is now possible to understand why Anomaly Detection is defined "model agnostic": in fact it does not rely on a specific theoretical model. This feature allows to scan a bigger range of events but it is less sensitive to any specific model. \\

% \begin{figure}[h!]
%     \centering
%     \includegraphics[width=2 in]{Images/Cap3/BumpHint.png}
%     \caption{Search for anomalies in particle physics. Figure from Ref. \cite{ML4ADHEP}}
%     \label{fig:bumphint}
% \end{figure}

\noindent To understand Anomaly Detection, we can use an example. Assuming we need to find a church in a town we do not know, without any technology. Instead of desperately walking around, we decide to climb an hill where we can scan the whole town with not so much effort. It is possible to identify the neighbourhood of the church or the "spatial region" where it is situated. We can also get some information we would not have got in the town, such as the colour of the roof or of the facade of the building we are looking for. Then we climb down the hill and go straight to the point we observed. Thanks to the details we obtained up there, we are able to find the church quite quickly and more easily. In other words, the church represents New Physics, climbing the hill is Anomaly Detection, climbing down is using normal data analysis techniques. It is quite powerful, isn't it? To apply Anomaly Detection it is necessary to hypothesize that anomalies are rare and they represent a small portion of a whole dataset. This assumption is realistic because it is what happens everyday in airports, finance and particle physics: normal events are way more common then deviations. In addition, if anomalies in Particle Physics were more frequent, we would have already observed them through common data analysis.

%%%%%%%%%%%%%%%%%%

\section{Deep Learning and Anomaly Detection}\label{sec:DLforAD}
While working on a Particle Physics problem, it is common to deal with a huge amount of data collected during collisions at the LHC. Machine Learning (ML) algorithms, which are able to learn from past by using vast datasets, can help in this sense. Among all categories of ML, Deep Learning (DL) is the best for our task \cite{LHCOlympics}. Based on a multi-layer Neural Network, DL can build complex data representation, which is optimal when dealing with high-dimensional dataset. Furthermore due to its implementation, that will be discussed in chapter \ref{chap:Model}, it can be exploited for model agnostic applications. \cite{ML4ADHEP}\\
ML and DL methods can be classified based on how the labels of data are used to train models. Main categories are: Supervised Learning, Unsupervised Learning and Semi-Supervised Learning.\cite{DLforAD}\\
In Supervised Learning, shown in figure \ref{fig:SUPtrain}, labeled data is used to train DL models and it can be implemented only with sufficient normal and anomalous data.

% \begin{figure}[h!]
%     \centering
%     \includegraphics[width=3.5 in]{Images/Cap3/SUPtraining.png}
%     \caption{Data used for Supervised Learning. Figure from Ref. \cite{DLforAD}}
%     \label{fig:SUPtrain}
% \end{figure}

\noindent Unsupervised Learning is based on not labeled data, as shown in figure \ref{fig:UNSUPtrain}, and models implementing it are trained to find a structure among the input features. Then model performance is tested using a mixed dataset containing both background and signal data.\\
Semi-Supervised Learning is a hybrid of previous techniques, implemented in the most varied ways.

% \begin{figure}[h!]
%     \centering
%     \includegraphics[width=2.7 in]{Images/Cap3/UNSUPtraining.png}
%     \caption{Data used for Unsupervised Learning. Figure from Ref. \cite{DLforAD}}
%     \label{fig:UNSUPtrain}
% \end{figure}

\noindent In this project we implemented an Unsupervised Anomaly Detection, because labels have not been used and model was trained only over background data.

%%%%%%%%%%%%%%%%%%%%%%
\section{How is it implemented?}\label{sec:ADimplementation}
The process of training a  DL model is divided into three main phases: training, validation and test loop. They are defined "loop" because in each of them a different section of the dataset is used and the code loops over it.\\
The most important stage is the training loop: in our case the chosen DL model is trained to reconstruct arrays containing information about particle jets, using only data with normal behaviour (background events) \cite{DLforAD}. In this phase, parameters of the model are adjusted in order to best fit to data. By doing so, DL model is able to recognise very well normal data and their structure. Our hope is that, once it faces an anomaly (signal event), it is able to identify it due to its difference from what it knows. The second step is the validation loop, where a valuation of the performance of the model is elaborated using the validation set. The valuation and the search for the best weights of the model are executed computing distance between output and input over training and validation set. When this distance is minimum within the validation loop, we save the model's weights.\\
The trained and tuned model is then tested, using a dataset containing unseen data, the test set, where there are both background and signal events. The purpose of this phase is to validate results previously obtained and get a preview of what will happen once the model will be running with completely new data. The task is completed computing the Anomaly Score (AS). In this thesis, we implement it as a distance between input and output, with different dimension: instead of calculating it over a batch of data, we do so for each event. In this way we are able to plot a reconstruction probability graph: we expect that the distributions of AS from background and signal events are different, because signal data contains anomalies which our model does not know. By means of these distributions we can identify a threshold of separation between them. For any event we pass to our trained model, we compute AS: if it higher (lower) than the chosen threshold, we are looking at a background (signal) event. Since AD is distinguishing signal events from background, we are basically executing a classification task with two classes. Implementation of Anomaly Detection is discussed more in depth in chapters \ref{chap:Model} and \ref{chap:Results}.

% \begin{figure}[h!]
%     \centering
%     \includegraphics[width=3 in]{Images/Cap3/ADtestloop.png}
%     \caption{Test loop of Anomaly Detection. Figure from Ref. \cite{DLforAD}}
%     \label{fig:RECOprob}
% \end{figure}


